\documentclass[12pt]{report}
\usepackage[T1]{fontenc}
\usepackage[brazilian]{babel}
\usepackage{csquotes}
\usepackage{makeidx}
\usepackage{glossaries}
\usepackage[
  perfil=projeto,
  impressao=frente-e-verso,
  citacoes=autor-data
]{abntex3}

\makeindex
\makeglossaries
\newglossaryentry{checksum}{
  name={soma de verificação},
  description={valor calculado para detectar alterações em dados}
}
\addbibresource{referencias-exemplo.bib}
\ABNTEXsetup{
  titulo={Modelo canônico completo de projeto de pesquisa},
  subtitulo={base auditável com todos os elementos implementados},
  autoria={Maria da Silva},
  instituicao={Universidade Exemplo},
  natureza={Projeto de pesquisa},
  objetivo={submetido ao Programa de Pós-Graduação em Documentação},
  area={Preservação digital},
  orientacao={Prof. Dr. João Souza},
  coorientacao={Profa. Dra. Ana Santos},
  local={Fortaleza},
  data={2026}
}
\ABNTEXabreviatura{cron.}{cronograma}
\ABNTEXsigla{PDF}{Portable Document Format}
\ABNTEXsimbolo{$t$}{Tempo de execução}

\begin{document}

% PARTE EXTERNA OPCIONAL
\ABNTEXprojetoexterna
\ABNTEXcapa
\ABNTEXlombada[Projeto 2026][Universidade Exemplo]

% ELEMENTOS PRÉ-TEXTUAIS
\ABNTEXprojetopretextual
\ABNTEXfolhaderosto
\ABNTEXlistadeilustracoes
\ABNTEXlistadetabelas
\ABNTEXlistadeabreviaturasesiglas
\ABNTEXlistadesimbolos
\ABNTEXsumario

% ELEMENTOS TEXTUAIS ESPECÍFICOS DO PERFIL
\ABNTEXprojetotextual
\ABNTEXprojetointroducao{%
  \index{projeto}
  \textbf{Tema:} preservação de documentos científicos digitais.

  \textbf{Problema:} quais procedimentos produzem evidências reprodutíveis de
  integridade e disponibilidade?

  \textbf{Hipótese:} formatos abertos, redundância e verificação automatizada
  reduzem perdas.

  \textbf{Objetivo geral:} comparar estratégias de preservação.

  \textbf{Objetivos específicos:} inventariar formatos, executar ensaios e
  documentar resultados.

  \textbf{Justificativa:} a continuidade do acesso sustenta pesquisa,
  transparência e memória institucional.
}
\ABNTEXprojetoreferencialteorico{%
  A normalização documental e a preservação digital delimitam o estudo
  \parencite{associacao2025}. Uma \gls{checksum} será usada nos ensaios.
  \begin{ABNTEXcitacaolonga}
  Texto demonstrativo para inspecionar o ambiente destinado às citações
  longas, sem transcrever conteúdo protegido.
  \end{ABNTEXcitacaolonga}
  \ABNTEXnota{Nota explicativa demonstrativa do projeto.}
}
\ABNTEXprojetometodologia{%
  A pesquisa será aplicada, documental e experimental. O universo será
  composto por arquivos de formatos distintos; a amostra será controlada.
  Os procedimentos incluem coleta, cálculo de somas, alteração deliberada,
  nova verificação e comparação dos resultados.

  \begin{figure}[ht]
    \centering
    \rule{0.55\linewidth}{2cm}
    \caption{Fluxo experimental proposto}
    \ABNTEXfonte{elaborada pela autora.}
  \end{figure}
}
\ABNTEXprojetorecursos{%
  Serão necessários repositório versionado, armazenamento redundante,
  ferramentas livres, estação de trabalho e horas da equipe.

  \begin{table}[ht]
    \centering
    \caption{Estimativa demonstrativa de recursos}
    \begin{tabular}{lr}
      Recurso & Quantidade\\
      Armazenamento & 2 unidades\\
      Estação de trabalho & 1 unidade
    \end{tabular}
    \ABNTEXfonte{elaborada pela autora.}
  \end{table}
}
\ABNTEXprojetocronograma{%
  \begin{center}
    \begin{tabular}{lll}
      Atividade & Início & Término\\
      Revisão documental & janeiro & março\\
      Experimentos & abril & julho\\
      Análise e redação & agosto & dezembro
    \end{tabular}
  \end{center}
}

% ELEMENTOS PÓS-TEXTUAIS
\ABNTEXprojetopostextual
\ABNTEXbibliografia
\ABNTEXglossario
\ABNTEXappendix{Instrumento de coleta proposto}
Formulário elaborado pela pesquisadora para registrar cada ensaio.
\ABNTEXannex{Especificação externa de formato}
Documento de origem externa que apoia o projeto.
\ABNTEXindice[Índice de assuntos]

\end{document}
